\chapter{基因检测能告诉我们什么}

\begin{kaichang}
想象你收到一个快递盒子。里面没有药品，没有仪器，只有一支塑料管、一小瓶保存液和一张说明书。按说明书的要求，你往管子里吐 2 毫升口水，盖紧，寄回。三个星期后，手机上弹出一份几十页的报告：你的酒精代谢能力“较弱”；你患某种疾病的风险“是平均水平的 1.3 倍”；你的祖源“约六成来自北方汉族”；你可能携带某个“意义未明”的基因变异。

现在先别管这份报告值多少钱，先猜一个问题：这一管口水，到底能“说出”多少关于你的真话？报告里的每一句话，可信度都一样吗？把直觉写下来。这一章结束时，你要学会用四个统计数字重新审读这类报告——而且你会发现，最需要警惕的往往不是检测技术，而是我们自己读百分比的方式。
\end{kaichang}

\section{口水里到底有什么}

先从看得见的说起。你吐进管子的那口液体，绝大部分是水，混着唾液里的酶和黏液。但对检测实验室来说，真正值钱的是漂在里面的几万个细胞——从你口腔内壁脱落的上皮细胞，还有少量白细胞。每个细胞核里，都装着一整套你的基因组：大约 30 亿个碱基对，用 A、T、C、G 四种“字母”写成的一部长书。第 66 章讲过这部书的结构，第 81 章讲过怎样逐字读它，这里不再重复，只问一个新问题：检测公司真的把 30 亿个字母全读了吗？

绝大多数消费级检测并没有。它们用的是一种叫“基因芯片”的取巧办法：不读全书，只抽查几十万个“常见拼写差异”的位置。这类位置有个专门名字——单核苷酸多态性，简称 SNP（读作“snip”）。意思是：在基因组的某个坐标上，人群里一部分人的字母是 A，另一部分人是 G，就像同一句话有人写“颜色”有人写“颜涩”。绝大多数 SNP 没有任何后果，只是人群里的自然拼写差异；但少数 SNP 恰好落在基因的开关上或蛋白质的关键部位，就会改变性状或风险。芯片像一张印好题目的答题卡，只问你那几十万个做过记号的坐标：“这个位置你是 A 还是 G？”便宜、快速，但注定看不到没做记号的地方。

医院里的检测则是另一种思路：高度怀疑某一种病时，把相关基因从头到尾完整测一遍，一个字母都不放过。两种做法没有谁高级谁低级——一个是广泛抽查，一个是定点精读，回答的问题本来就不一样。记住这个区别，因为后面所有的“报告可信度”讨论，都要先问：它读的是什么？

\section{同一管口水，三种分量完全不同的答案}

基因检测报告最容易骗人的地方，是把分量天差地别的结论印在同一个页面上、用同一种字体。我们按“答案的硬度”把它们分成三类。

\textbf{第一类：单基因病——接近“是或否”的答案。}有些疾病几乎由一个基因的拼写错误决定：囊性纤维化（一个编码氯离子通道的基因坏了，第 15 章讲过的离子运输在这里出问题）、镰刀形贫血（血红蛋白上一个氨基酸的替换，第 72 章讲过的错义突变）、亨廷顿病（本章后面会细讲）。对这类病，只要检测技术本身可靠，找到致病变异就基本等于宣判，没找到就基本等于排除。这类答案最“硬”，但请注意限定词：“几乎”和“基本”——即便是单基因病，也存在携带着变异却始终不发病的例外（遗传学管这叫“外显不全”），还存在一种最磨人的结果：“意义未明变异”，也就是发现了一处拼写变化，但全世界的科学家都还不能确定它有没有害处。

\textbf{第二类：复杂疾病风险——答案只是一张“概率微调单”。}高血压、2 型糖尿病、冠心病、大多数癌症和抑郁症，都不是一个基因说了算的。它们背后往往有几百甚至几千个位点，每个位点只把风险往上或往下推一点点，再加上饮食、运动、睡眠、污染、运气。基因检测对这类病给出的是“多基因风险评分”：把所有已知位点的微小影响加起来，告诉你落在人群分布的哪个百分位。这是统计学意义上的真实信息，但它移动的是概率，不是命运——这一点我们马上要用数字严格地算给你看。

\textbf{第三类：药物基因组学——当下最“立竿见影”的一类。}基因不仅决定你会得什么病，还决定你的身体怎样处理药物。最贴近生活的例子是喝酒脸红：东亚人群中相当比例的人携带一种乙醛脱氢酶（\chem{ALDH2}）的变异版本，酒精代谢的中间产物乙醛分解得慢，堆积在血液里，于是脸红、心跳、难受。这不只是“能不能喝”的社交问题——乙醛是一类明确的致癌物，带着这种变异还长期饮酒，食管癌风险显著升高。更严肃的例子在心内科：抗血栓药氯吡格雷本身没有药效，必须靠肝脏里的一种酶（\chem{CYP2C19}）把它转化成活性形式；而人群中有一部分人天生这种酶活性低，常规剂量对他们几乎无效，装完心脏支架照样可能再次血栓。现在不少医院在开这种药之前会先查基因，这就是药物基因组学的日常应用：不预测疾病，只校准用药。

同一管口水，给出的是“是或否”、“概率微调”和“用药说明书”三种完全不同的东西。混在一起读，是误解的第一个源头。

\section{四个数字：审判任何一份检测报告}

现在进入本章最硬的一节。不管是基因检测、体检筛查还是验孕棒，任何一个“检测”的性能都由四个数字描述。这四个数字你一辈子都会用到，值得慢慢读。

\begin{qiaoliang}
把所有人按“真实状态”分成两列：真的有病、真的没病；再按“检测结果”分成两行：阳性、阴性。得到一张四个格子的表：

\begin{center}
\begin{tabular}{lcc}
\toprule
 & 真实：有病 & 真实：没病 \\
\midrule
检测：阳性 & 真阳性 & 假阳性 \\
检测：阴性 & 假阴性 & 真阴性 \\
\bottomrule
\end{tabular}
\end{center}

\textbf{敏感度}：有病的人里，检测能抓出多少比例？真阳性除以所有病人。敏感度 99\%，意思是 100 个真病人里能逮住 99 个，漏掉 1 个。

\textbf{特异度}：没病的人里，检测能放对多少比例？真阴性除以所有健康人。特异度 95\%，意思是 100 个健康人里会冤枉 5 个，制造 5 个假阳性。

大多数人以为这两个数字高，检测就可信。下面这个计算会证明：\textbf{还差一个隐藏的主角——基础概率}，也就是这种病在人群里的患病率。

假设某种病的患病率是 1/1000（千分之一），某检测敏感度 99\%、特异度 95\%——看起来相当优秀。现在让 10 万人接受检测：

\begin{itemize}[itemsep=2pt]
\item 10 万人里大约有 100 个真病人（$100000 \times 1/1000$）。敏感度 99\%，逮住约 99 个真阳性。
\item 剩下 99900 个健康人，特异度 95\% 意味着 5\% 被冤枉：约 4995 个假阳性。
\item 阳性结果总数：$99 + 4995 = 5094$。其中真病人只有 99 个。
\item 所以，拿到阳性结果的人里真正患病的比例——这叫\textbf{阳性预测值}——是 $99/5094 \approx 1.9\%$。
\end{itemize}

一个“99\% 敏感”的优秀检测，阳性结果竟然约 98\% 是误报。这不是检测的错，是算术的错：罕见病的真病人太少，健康人太多，哪怕冤枉健康人的比例很小，冤枉出来的绝对人数也会淹没真病人。\textbf{基础概率越低，阳性结果越可疑}——这是本章最重要的一句话。
\end{qiaoliang}

这就是为什么正规医学把“筛查”和“确诊”分成两步：先用便宜、敏感度高的检测把可疑的人筛出来（宁可错杀，不可放过），再用昂贵、更精确的方法（比如活检、基因全长测序）复核。两步走，就是用第二次检测把基础概率抬高——进入第二步的人已经经过第一轮筛选，患病率不再是 1/1000，而是高得多，同一个检测的阳性预测值立刻变得好看。这个思路，等下你将在实验里亲手摸到。

\begin{tiaozhan}
上面的计算其实是一条定理的特例：贝叶斯定理。用符号写出来就是
\[P(\text{有病} \mid \text{阳性}) = \frac{P(\text{阳性} \mid \text{有病}) \times P(\text{有病})}{P(\text{阳性})}\]
其中 $P(\text{有病})$ 就是基础概率，$P(\text{阳性} \mid \text{有病})$ 就是敏感度。分母 $P(\text{阳性})$ 是“所有阳性结果”的总概率，要把真阳性和假阳性加在一起。这条定理说的是：\textbf{新证据（阳性结果）能把你的判断更新多少，取决于你原来的判断（基础概率）有多强}。证据不变，起点不同，结论就不同。跳过这段符号不影响主线，但如果你记住“结论 = 证据 × 起点”，就已经掌握了本章的数学灵魂。
\end{tiaozhan}

\section{关联不等于因果，风险不等于诊断}

解决了“检测准不准”，还剩两个更隐蔽的陷阱。

第一个陷阱：\textbf{关联不等于因果}。复杂疾病的风险位点是怎样找到的？方法是“全基因组关联研究”（GWAS）：招募几万名病人和几万名健康人，比较几十万个 SNP 在两组里的出现频率，凡是在病人组明显更常见的位点，就记为“与该疾病关联”。这个方法很强大，但它的词典里只有“关联”这个词。夏天冰淇淋销量和溺水人数高度关联——两者都是“天气热”的产物，谁也不是谁的原因。同理，一个关联位点可能只是恰好住在真正致病因子的隔壁——染色体上相邻的片段往往整体遗传（第 64 章讲过的连锁），所以它是个路标，不一定是肇事者；它也可能标记的是人群差异：比如某个位点在北欧人群里常见，而这个人群恰好平均身高更高，粗心的分析就会把“长高基因”的头衔错发给一个无辜的路标。从关联到因果，还需要细胞实验、动物实验等一整套后续工作去验证机制（第 81 章的伏笔在这里回收：能读序列，不等于懂了功能）。

第二个陷阱：\textbf{风险不等于诊断}。报告里最常见的写法是“你的风险是平均水平的 1.3 倍”。听起来吓人，但请养成一个习惯：立刻追问“平均水平是多少”。如果这种病的平均终身风险是 1\%，1.3 倍就是 1.3\%——每千人里 10 个变成 13 个。相对风险翻一倍，绝对风险可能只挪动了小数点。反过来，名气很大的基因也一样受这条规则约束：携带 \chem{BRCA1} 基因某些突变的女性，一生中患乳腺癌的风险确实大幅升高（从普通人群的约十分之一升至约一半甚至更高），这是有坚实证据的风险信息，值得认真对待——但“约一半”仍然不是“诊断”，携带者中仍有相当比例终身不发病；而没有这种突变，也不等于上了保险，因为绝大多数乳腺癌恰恰发生在没有已知突变的普通人群里。风险信息的价值在于帮你和医生决定“要不要更早、更勤地筛查”，而不在于宣判。

\begin{wuqu}
“检测结果是阳性，说明我有病（或者迟早要得病）。”——这是关于检测流传最广的误解，本章的前两节就是它的双重反例。第一重：基础概率低时，阳性结果大部分是假阳性，10 万人例子里 98\% 的阳性者其实健康。第二重：即使是针对真突变、技术上零误差的检测，对复杂疾病它给出的也只是概率的挪动。反过来说也成立：阴性结果同样不等于高枕无忧——芯片只查了做了记号的位点，没查到的变异、没发现的位点、环境和运气，都还在牌桌上。检测报告是概率世界的对话，不是判决世界的文书。
\end{wuqu}

\section{祖源报告：一张会自己变化的饼图}

消费级报告里最好看的部分往往是那张五颜六色的祖源饼图：百分之多少来自北方汉族，百分之多少来自东南亚，说不定还有 2\% 的“尼安德特人成分”。它也是全书最需要“打引号”的部分，因为它的生成方式和前面所有检测都不同：\textbf{它几乎不测任何“事实”，只做统计匹配}。

原理是这样的：检测公司先建立一个“参考人群库”——收集成千上万份已知祖源的 DNA 样本，把每个群体的 SNP 模式统计出来。然后把你的 SNP 模式和库里的各个人群逐一比对，看你的每一段序列“最像”库里的哪个群体，最后把相似度折算成百分比。看出来了吗？结论完全依赖两个可以人为改变的东西：\textbf{参考库里收了谁}，和\textbf{比对用什么统计模型}。库里如果几乎没有中亚样本，你祖上明明有中亚血统，报告也只能把你硬塞进“最接近”的类别。所以同一个人把样本寄给两家公司，会得到两张不同的饼图；甚至同一家公司更新数据库之后，你几年前的报告会自己“变脸”——你的 DNA 一个字母都没变，变的是参照系。

更根本的边界是：这张饼图说的是“你的某些片段和\textbf{今天}哪群人最相似”，而不是“你的祖先几千年前生活在哪块土地上”。今天的人群是几千年迁徙、通婚、分分合合之后的快照，拿快照反推历史，天然带着模糊。祖源检测是一种有趣的人口学游戏，偶尔也能帮被收养者找到血亲，但它不是血统证明书，更不该成为“我是谁”的身份裁判。

\section{一个家族、一个湖和一段被数出来的重复}

\begin{zhengju}
基因检测史上最著名的证据故事，开始于一个家族的悲剧。美国神经学家南希·韦克斯勒的母亲和三位舅姨都死于亨廷顿病——一种中年起病的遗传病：先是手脚不自主地舞动，接着认知衰退，最后在发病十几年后离世。如果父亲或母亲患病，每个子女都有 50\% 的概率继承它。南希自己就是一个“抛在半空的硬币”。

1979 年起，韦克斯勒带队反复前往委内瑞拉的马拉开波湖畔。那里的几个村庄有全世界最大的亨廷顿病家系——一万多人的大家族，许多人的血缘可以追溯回两百年前的一位共同祖先。研究团队挨家挨户绘制家谱、记录谁发病谁健康、采集血样。问题很明确：在连基因长什么样都不知道的情况下，能不能先在染色体上找到它的“门牌号”？

方法来自第 64 章的连锁思想：既然致病变异和它的“邻居”在染色体上常常整体遗传，那就找一种能看出来的邻居——当时用的是 DNA 上天然存在的“限制性内切酶切点差异”（第 82 章会讲这类酶）：如果某个切点变异在患者家族里总是和疾病一起出现，致病变异就住在它附近。1983 年，团队真的找到了这样一个标记，把致病基因定位到第 4 号染色体——尽管当时没人知道那是什么基因、什么蛋白质。这一下，“症状前检测”成为可能：年轻人可以在发病前几十年知道自己是否携带它。但请注意这个故事的诚实部分：连锁标记不是基因本身，染色体在遗传时会发生交换（重组），偶尔会把标记和致病变异“拆散”，所以早期检测存在小比例的假阳性和假阴性——人们不得不联合多个标记、要求多位家族成员参与，才能把错误率压低。又过了整整十年，1993 年，国际合作团队才终于抓到基因本身：一段 CAG 三字母重复，正常人重复十几次到三十几次，患者重复 40 次以上，而且重复次数越多，发病往往越早。

这个故事留给我们的不只是“找到了”三个字。第一，它展示了一条完整的证据链：从家谱（现象）到连锁标记（统计关联）再到基因序列和重复次数（分子机制），每一步都在回答“会不会有别的解释”。第二，它逼出了一个至今没有标准答案的伦理问题：一种无法治愈的病，检测能告诉你未来，却不能改变未来——你要不要提前知道？结果出人意料地分化：许多真正处于风险中的人，权衡之后选择\textbf{不检测}。正是围绕亨廷顿病检测，医学界建立起今天通用的规矩：检测前必须有遗传咨询、必须本人知情同意、必须保护结果不被滥用。科学给出了“能测”的能力，社会必须为“要不要测、谁能知道”立规矩。
\end{zhengju}

\section{动手试一试：亲手制造一批“假阳性”}

\begin{shiyan}
\textbf{材料}：100 张大小相同的纸条、一个不透明的袋子、两个小盒子（或两个信封）、笔。不需要任何生物材料，也不要使用任何人的真实样本——这只是一场概率游戏。

\textbf{步骤}：
\begin{enumerate}[itemsep=1pt]
\item 做“人群”：在 2 张纸条上写“病”，98 张上写“健康”，折好放进袋子混匀。你造出了一个患病率 2\% 的人群。
\item 做“检测仪”：第一个盒子里放 10 张卡，9 张写“阳性”、1 张写“阴性”——这是给“病人”用的检测（敏感度 90\%）；第二个盒子里 1 张写“阳性”、9 张写“阴性”——这是给“健康人”用的检测（特异度 90\%）。
\item 每抽一张人群纸条，根据上面写的是“病”还是“健康”，去对应的盒子里抽一张结果卡，看一眼，记下结果，然后把结果卡放回盒子摇匀。纸条也放回袋子——这样每个人的检测互不干扰。
\item 重复至少 50 次（可以几个同学分工，把数据合并），只统计一件事：\textbf{所有“阳性”结果里，有多少次对应的纸条真是“病”？}
\end{enumerate}

\textbf{观察点}：理论上，100 次检测里约 2 次碰到真病人、其中约 1.8 次报阳性；约 98 次碰到健康人、其中约 9.8 次报假阳性。所以阳性里约 12 次中，真病人只占约 2 次——阳性预测值只有约 15\%，尽管你的检测仪“敏感度和特异度都高达 90\%”。

\textbf{变式}：把“病”纸条改成 20 张（患病率 20\%），重做一遍。同一个检测仪，阳性预测值会蹿升到约 69\%。亲手做一次，你对“基础概率”四个字的理解会超过读十遍定义。

\textbf{安全提示}：本实验只用纸和笔。真正的基因检测涉及他人不可更改的遗传信息，绝不能在未经本人同意的情况下采集、检测任何人的生物样本。
\end{shiyan}

\section{你的基因数据，不只是你的}

最后这一段不讲科学，讲边界——因为基因检测的特殊之处，恰恰在于它测的东西特殊。

第一，\textbf{DNA 无法挂失}。密码泄露了可以改，指纹泄露了改不了，而基因组泄露了，你连“换一根手指”的选择都没有——它终身有效，还包含你尚未发病时的未来信息。

第二，\textbf{它天然是“共享文件”}。你的一半基因组来自父亲、一半来自母亲，你的子女各携带你的一半。你一个人上传数据，等于部分公开了父母、兄弟姐妹、子女的信息——而他们未必同意过。2018 年，美国警方正是利用这一特征破获了悬置四十年的“金州杀手”案：凶手本人从未上传过任何样本，是他的远亲把数据传进了公开的家谱数据库，警方顺着亲缘网络锁定了他。正义得到伸张，但争议至今没有平息：每一个上传者，都在无意中把几百个远亲变成了可被检索的对象。

第三，\textbf{歧视的风险真实存在}。雇主如果知道求职者携带某种高风险基因，会不会拒绝录用？保险公司会不会拒保或加价？一些国家已经立法禁止基因歧视（美国在 2008 年通过了《遗传信息无歧视法案》），中国法律也对个人信息和人类遗传资源有严格保护，但立法永远跑在技术的后面，自我保护意识不能缺席。

第四，\textbf{知情同意不是走过场}。正规的临床基因检测之前，医生或遗传咨询师必须告诉你：测什么、不测什么、可能意外发现什么（比如“检出”父子之间没有血缘关系——这类“附带发现”远比一般人想象的常见）、数据会被保存多久、会不会用于研究。任何不给你这些信息的检测，都不值得信任。

而把上面所有原则收拢成一句最实用的忠告，就是本章大纲要求的最后一句话：\textbf{不要依据任何一份消费级检测报告，自行改变你的医疗方案}——不要自行停药、加药，不要自行决定做或不做手术。消费级检测的芯片覆盖有限、阳性预测值受基础概率摆布、变异解读常常“意义未明”。报告里真正让你不安的内容，正确去处是医院的临床级复检和遗传咨询门诊，而不是搜索引擎和朋友圈。

\section{回到那管口水}

现在兑现开场的承诺，重新读那三句话。

“酒精代谢能力较弱”——这属于单一位点、机制清楚的判断（乙醛脱氢酶的活性差异），可信度很高，而且是罕见的“可以直接指导行为”的结果：脸红不是能练出来的酒量，是身体在报警。

“患病风险是平均水平的 1.3 倍”——这是多基因评分，建立在人群关联之上。它没说你有病，甚至没说你多半会得病；它说的是在一个基础概率可能只有百分之几的病上，你的数字略微抬高了一点。要不要据此改变生活方式？注意：无论你得分高低，不熬夜、多运动、不吸烟都是对的——这类建议的正确性根本不需要基因检测来证明。

“约六成来自北方汉族”——这是统计模型的输出，参照系是某家公司某个版本的参考库。它描述了相似性，不证明来源，明年更新数据库它自己就会改口。

一管口水确实装着你的全部遗传信息，但“装着信息”和“能回答问题”之间，隔着检测技术、统计模型、基础概率和解读者的诚实。基因检测是一台强大的仪器——强大到能预言一个家族的命运；同时也普通到必须接受四条统计规律的审判。学会审读它，你就不会被它吓到，也不会被卖它的人骗到。

\section{往更大的尺度走一步}

这一章里有一个反复出现的潜台词，值得你停下来注意：基因的意义，几乎从来不由基因自己说了算。同样的乙醛脱氢酶变异，在不喝酒的文化里无关痛痒，在劝酒文化里就是致癌风险的放大器；同样的 BRCA1 突变，在有定期乳腺筛查的医疗体系里和在没有筛查的地方，通向完全不同的结局。基因总是在和环境、群体、文化“合写”结局——下一章我们要把这个视角推到最大尺度：离开个体，去看整个生态系统怎样作为一个化学与信息的网络运行（第 86 章）。而在基因组内部，人类读懂的部分其实也只有一小角：那些不编码蛋白质、占基因组绝大部分的序列还有多少未知功能，是最后一章（第 87 章）要直面的未解之谜之一。

\section*{练一练}

\begin{sikaoti}
\item 一种病患病率为 1/100（百分之一），某检测敏感度 90\%、特异度 90\%。仿照本章的 10 万人算法，画出四格表，算出：检测呈阳性的人里，真正患病的比例是多少？再和本章 1/1000 的例子对比，用自己的话说出差别的原因。
\item 某报告称“携带某变异者患某病风险是普通人的 2 倍”。已知普通人患此病的终身风险是 0.5\%。请算出携带者的绝对风险，并回答：1000 个携带者里，预计有多少人终身不会得这种病？
\item 画一张信息流图：唾液 $\rightarrow$ 细胞 $\rightarrow$ DNA $\rightarrow$ 基因芯片 $\rightarrow$ SNP 数据 $\rightarrow$ 统计模型 $\rightarrow$ 报告上的百分比。在每一级箭头上各写一种“可能引入错误或不确定性”的方式（提示：至少六级，至少六种来源）。
\item 小明的祖源饼图从“70\% 北方汉族”变成了“62\% 北方汉族、8\% 东北亚”，而他的 DNA 没有任何变化。用“参考人群库”和“统计模型”解释这种变化，并说明它为什么不能作为血统证明。
\item 小华私下买了基因检测盒，用父亲用过的牙刷做了检测，并把父亲“高风险”的结果发到了家庭群里。请从本章讲的至少三条原则出发，指出这个行为错在哪里。
\item 假设未来某种多基因风险评分能把 2 型糖尿病高危人群识别得很准，但该病目前无法根治、只能靠生活方式干预延缓。仿照亨廷顿病的伦理讨论，说说你会支持还是反对把这项检测普及给所有中学生，给出至少两个理由。
\end{sikaoti}
